Protein language models and related variant effect predictors (VEPs) are increasingly used to classify missense variants, yet their reliability is typically summarized by benchmark-wide averages, and they are applied to proteins evolutionarily distant from those used for calibration: a family-level distribution shift. Standard split conformal prediction assumes exchangeable variants, and existing VEP calibration schemes condition on residue, gene or domain attributes rather than on homology support.
We present ProtCOVER, a conformal procedure that treats homology families, rather than individual variants, as the exchangeable unit. ProtCOVER (i) calibrates at the family level over label-free homology clusters, yielding coverage for a new family under family exchangeability; (ii) adapts to descriptor shift by calibrating conditionally on an estimated family-level density ratio of label-free descriptors (sequence similarity, alignment depth and embedding distance), so that family-weighted coverage holds exactly under family exchangeability, under the estimated shift and under every reweighting in their span, and within a total-variation bound otherwise; and (iii) computes each query's effective calibration support and abstains when it is insufficient.
On ProteinGym clinical substitutions with five zero-shot predictors, Mondrian split conformal covers only \sOneBtwoPath{} of pathogenic variants per family at a 0.90 target, even without shift. Under descriptor shift, ProtCOVER's mean worse-class shortfall below the target is \shortCrTwoMean{}, against \shortGwcpTwoMean{} for group-weighted and \shortHcpTwoMean{} for hierarchical conformal prediction; on orphan families and held-out embedding regions, which do not shift its descriptors, it matches group weighting (\shortCrIndMean{} against \shortGwcpIndMean{}). The ordering holds with deployment-realistic, alignment-free features. ProtCOVER's sets are larger under shift but not at matched coverage. A shift in predictor accuracy that the descriptors cannot see defeats every method, and support-based abstention is no better than random abstention. In a pre-registered evaluation on ClinVar variants added after the benchmark's snapshot, every shift-aware method met the target. Code is publicly available.
